% TOP_PROBLEM: {"id": 85, "title": "Hilbert's 7th Problem: Transcendence of Certain Numbers", "category_id": 1, "category": {"id": 1, "name": "number_theory", "display_name": "Number Theory", "description": "Properties of integers, prime numbers, Diophantine equations.", "slug": "number-theory", "order_index": 1, "created_at": "2026-07-31T15:26:25.670Z"}, "status": "open"}
% FIRST_POSTED: null
% ENTRY_KIND: statement_only
% Imported from: batch02.tex; UnsolvedMath ID 85
\documentclass[11pt]{article}
\usepackage[margin=25mm]{geometry}
\usepackage{fontspec}
\setmainfont{Latin Modern Roman}
\usepackage{amsmath,amssymb,mathtools}
\usepackage{unicode-math}
\setmathfont{Latin Modern Math}
\usepackage{xurl,hyperref}
\hypersetup{colorlinks=true,urlcolor=blue}
\setcounter{secnumdepth}{0}
\setlength{\parindent}{0pt}
\setlength{\parskip}{5pt}
\setlength{\emergencystretch}{3em}
\newcommand{\sref}[2]{\hyperlink{source:#1:#2}{[S#2]}}
\newcommand{\cataloguescope}[1]{\paragraph{Recorded target.}#1}
\begin{document}
% UnsolvedMath ID 85; source catalogue code HIL-007
\setcounter{section}{84}
\section{Hilbert's 7th Problem: Transcendence of Certain Numbers}\label{85}
{\small\sffamily UnsolvedMath ID 85 · Number Theory}
\subsection{Definitions and mathematical statement}

\paragraph{Shared notation.}$\mathbb N=\{1,2,\ldots\}$, and $\mathbb Z,\mathbb Q,\mathbb R,\mathbb C$ denote the integers, rationals, reals and complex numbers. Any other conventions are specified locally.


A complex number is \emph{algebraic} if it is a zero of a nonzero polynomial in $\mathbb Q[X]$; otherwise it is \emph{transcendental}. Let $\overline{\mathbb Q}\subset\mathbb C$ be the field of algebraic numbers. Here \emph{irrational} means not in $\mathbb Q$. For $\alpha\ne0$, the values of the possibly multivalued power $\alpha^\beta$ are $\exp(\beta L)$ as $L\in\mathbb C$ ranges over the logarithms satisfying $\exp(L)=\alpha$.
\paragraph{Statement recorded in the supplied source.}
For every $\alpha,\beta\in\overline{\mathbb Q}\setminus\mathbb Q$ and every $L\in\mathbb C$ such that $\exp(L)=\alpha$, is $\exp(\beta L)$ transcendental? Equivalently, is every value of $\alpha^\beta$ transcendental when both $\alpha$ and $\beta$ are algebraic and irrational? \sref{85}{1}\sref{85}{2}
This historical question has an affirmative answer by the Gelfond--Schneider theorem (1934), which also permits rational algebraic bases other than $0$ and $1$. The theorem applies to every fixed choice of logarithm. \sref{85}{2}
\subsection{Short English statement}
Must raising an algebraic irrational number to an algebraic irrational power always give a transcendental number, for every complex value of the power? This historical question is answered affirmatively by the Gelfond--Schneider theorem.
\subsection{Sources}
\begin{description}
\item[\hypertarget{source:85:1}{S1}] ULAM.AI and the UnsolvedMath Contributors, \emph{UnsolvedMath: A Curated Collection of Open Mathematics Problems}, record 85 (HIL-007). CC BY 4.0. \url{https://huggingface.co/datasets/ulamai/UnsolvedMath}.
\item[\hypertarget{source:85:2}{S2}] Michail Karatarakis and Freek Wiedijk, \emph{A formalization of the Gelfond--Schneider theorem}, arXiv:2603.24823v1, 25 March 2026, section 1, pp. 1--2, Theorem 3.1 and section 3.1, pp. 5--6 (including arbitrary logarithm values). \url{https://arxiv.org/pdf/2603.24823v1}.
\end{description}
\label{end:85}
\end{document}
